\subsection{DET-12: the direct signed Gram kernel and its valid positive geometry}
This section settles a specific proposed passage from the actual finite scalar to
positive operators. It also records the positive determinant-line and projection
factorizations that the scalar does possess. The obstruction concerns the direct
scalar; suitable transforms are a separate construction.

\paragraph{Source convention and exact input.}
The source \path{Formalism/FormalDegreeBessel/StatementG.lean}, line 9, uses
\(\operatorname{herm}(v,w)=w^*Hv\). In the paper's global conjugate-first convention
we write \(h(u,v)=u^*Hv\), so \(\operatorname{herm}(v,w)=h(w,v)\).
Lines 20--24 define \(H=H_{k,r}\), of signature \((r+k,r+1)\).
Lines 28--31 give the middle-support condition and \(T(z)=-h(z,z)\).
The latter is real because \(H\) is Hermitian. Lines 61--71 define the flag
vectors, cross-pairing matrix \(\widehat Z\), and leading minors \(D_l\);
lines 85--99 define
\[
 G_f=|D_{r+1}|^2-T(z)T(z')|D_r|^2
\]
and the Gram matrix of
\[
 e_0,\ldots,e_{r-1},z,ge_0,\ldots,ge_{r-1},gz'.
\]
We retain the exact hypotheses: middle-supported \(z,z'\), \(T(z),T(z')>0\),
and \(g^*Hg=H\). The proposition definitions \texttt{GGram} and
\texttt{GNonneg} are at lines 165--172; actual theorem exports
\texttt{gGram} and \texttt{gNonneg} are in
\path{Formalism/FormalDegreeBessel/FinalG.lean}, lines 21--22.
The balanced-inertia mechanism was independently proved in \fid{DET-10}.

\paragraph{DET-12(a): no direct positive Gram kernel.}
Set \(r=0\), \(k=2\), and \(H=\operatorname{diag}(1,1,-1)\).
On the negative unit vectors define
\[
 K(z,z')=|h(z,z')|^2-1=G_f(2,0,z,z',I).
\]
This is not a positive semidefinite kernel, hence is not a Hilbert Gram
kernel or a pointwise limit of positive semidefinite kernels.
The same conclusion holds for the direct normalized scalar
\((|h(z,z')|^2-T(z)T(z'))/(T(z)T(z'))\).
\begin{proof}
For \(t\in\mathbb R\), put \(z(t)=(\sinh t,0,\cosh t)^T\).
The support condition is vacuous for \(r=0\),
\(h(z(t),z(t))=-1\), and \(I^*HI=H\). Thus these are actual
source-admissible data. Direct calculation gives
\[
 h(z(s),z(t))=-\cosh(t-s),\qquad
 K(z(s),z(t))=\sinh^2(t-s).
\]
For \(t\ne0\), the two-point matrix at \(z(0),z(t)\) is
\[
 \begin{pmatrix}0&c\\c&0\end{pmatrix},\qquad c=\sinh^2t>0.
\]
Its quadratic form at \((1,-1)\) is \(-2c<0\).
A Hilbert Gram kernel has nonnegative quadratic forms, so no such
realization exists. Equivalently, the zero diagonal would force every
feature vector to be zero while the off-diagonal is nonzero.
For a pointwise limit of positive kernels, each fixed finite quadratic
form is the limit of nonnegative numbers and is nonnegative.
The displayed matrix rules out that possibility, for sequences or nets.
The normalized scalar restricts to this same kernel on the negative
unit vectors, proving its assertion as well.
\end{proof}

\paragraph{The group and positive trace-coefficient versions.}
Let \(z_0=(0,0,1)^T\) and
\[
 F(g)=G_f(2,0,z_0,z_0,g)=|h(z_0,gz_0)|^2-1,\qquad g^*Hg=H.
\]
Then \(F\) is not of positive type, is not a pointwise limit of
positive-type functions, and cannot be \(\Tr(A\pi(g))\) for a unitary
representation \(\pi\) and positive trace-class operator \(A\).
\begin{proof}
The matrices
\[
 g(t)=\begin{pmatrix}
 \cosh t&0&\sinh t\\0&1&0\\\sinh t&0&\cosh t
 \end{pmatrix}
\]
satisfy \(g(t)^*Hg(t)=H\), \(g(t)z_0=z(t)\), and
\(g(s)g(t)=g(s+t)\), by matrix multiplication and the hyperbolic
addition identities. Hence \(F(e)=0\) and \(F(g(t))=\sinh^2t\).
The positive-type matrix \(F(g_i^{-1}g_j)\) for \(g_1=e,g_2=g(t)\)
is the indefinite matrix above, since \(g(t)^{-1}=g(-t)\).
The pointwise-limit argument is unchanged.

Every proposed positive trace coefficient would be of positive type:
for \(X=\sum_j c_j\pi(g_j)\),
\[
 \sum_{i,j}\overline{c_i}c_j\Tr(A\pi(g_i^{-1}g_j))
 =\Tr(AX^*X)=\|XA^{1/2}\|_{\rm HS}^2\ge0.
\]
Here \(A^{1/2}\) is Hilbert--Schmidt and the cyclic trace identity
applies to the bounded and Hilbert--Schmidt factors.
Alternatively, \(F(e)=\Tr A=0\) gives
\(\|A^{1/2}\|_{\rm HS}=0\), hence \(A=0\), contradicting
\(F(g(t))>0\).
\end{proof}

\paragraph{The direct scalar is not a squared Hilbert distance.}
For \(a>0\), take \(t=(-a,0,a)\) and real coefficients \(c=(1,-2,1)\).
Their sum is zero, but
\[
 \sum_{i,j}c_ic_jK(z(t_i),z(t_j))
 =2\sinh^2(2a)-8\sinh^2a=8\sinh^4a>0.
\]
The last equality follows from \(\sinh(2a)=2\sinh a\cosh a\) and
\(\cosh^2a-1=\sinh^2a\). Thus conditional negative semidefiniteness
fails. Indeed, a squared Hilbert distance would satisfy, by expansion
and \(\sum_i c_i=0\),
\[
 \sum_{i,j}c_ic_j\|\Phi_i-\Phi_j\|^2
 =-2\left\|\sum_i c_i\Phi_i\right\|^2\le0.
\]
This proves the stated distance obstruction completely. These
refutations concern the direct scalar and the specified normalization,
not an inverse, exponential, integral transform, added or subtracted
term, or another modified family.

\paragraph{DET-12(b): positive determinant lines, indefinite ambient exterior form.}
Write \(q=r+1\), and let \(X\) contain the \(2q\) source Gram vectors.
The induced exterior Hermitian form is specified by
\[
 h^{[2q]}(x_1\wedge\cdots\wedge x_{2q},
          y_1\wedge\cdots\wedge y_{2q})
 =\det[h(x_i,y_j)]_{i,j}.
\]
For the column wedge \(\omega\) of \(X\), the source Gram identity gives
\[
 G_f=(-1)^q h^{[2q]}(\omega,\omega).
\]
If \(X^*HX\) is nonsingular, the signed form on the determinant line
\(\Lambda^{2q}(\operatorname{im}X)\) is positive definite. For \(k\ge2\),
however, that signed form is indefinite on the full ambient exterior
power.
\begin{proof}
By \fid{DET-10}, \(V=\operatorname{im}X\) has signature \((q,q)\).
In an \(h\)-orthonormal basis of \(V\), the exterior norm of its top
wedge has sign \((-1)^q\); multiplication by the same sign gives
value \(1\). Every other nonzero element of this one-dimensional line
is a nonzero scalar multiple, so the signed form is positive definite.

For the ambient signature \((p,q)=(r+k,r+1)\), a basis wedge
containing \(q\) positive and \(q\) negative directions has signed
value \(+1\). When \(k\ge2\), there are also \(q+1\) positive directions,
so a wedge containing \(q+1\) positive and \(q-1\) negative directions
has signed value \(-1\). This proves ambient indefiniteness.
When \(k=1\), the ambient dimension is \(2q\), so its top exterior
space is the positive line just described. When \(k=0\), the exterior
degree exceeds the ambient dimension, and that exterior space is zero.
For a degenerate restricted form, the displayed wedge has value zero,
so it does not define a positive definite metric on that line.
\end{proof}
The source compound matrices and wedge coordinates are
\path{Development/AbelianVarietyTheory/Definitions/Jets/SumIoPrAlt.lean},
lines 104--110. Its \texttt{TensH}, lines 88--94, still carries
signs \(\epsilon_I\in\{1,-1\}\). The determinant-line observation
does not remove those signs from the actual period expansion.

\paragraph{DET-12(c): explicit Lorentz-projection factorization.}
Let \(E=[e_0,\ldots,e_{r-1}]\), \(F=gE\), \(y=gz'\), and
\[
 C=E^*HF,\quad b=E^*Hy,\quad c=z^*HF,\quad d=z^*Hy.
\]
Then \(\widehat Z=\left(\begin{smallmatrix}C&b\\c&d\end{smallmatrix}\right)\)
and \(D_r=\det C\). On the chart \(\det C\ne0\), define
\[
 U=\operatorname{im}E+\operatorname{im}F,\quad
 x=z-E(C^*)^{-1}c^*,\quad y_0=y-FC^{-1}b,\quad
 s=d-cC^{-1}b.
\]
With \(T=T(z)>0\), the space
\(V_+=U^\perp\cap x^\perp\) is positive definite of dimension \(k\),
the vector \(p=y_0+(s/T)x\) belongs to it, and
\[
 \boxed{\quad G_f=|\det C|^2T\,h(p,p)
 =\bigl\|\sqrt T\,(\det C)p\bigr\|_{V_+}^{\,2}.\quad}
\]
\begin{proof}
The source support, isotropy and unitary conditions give
\begin{align*}
 E^*HE&=F^*HF=0,\qquad E^*Hz=0,\qquad F^*Hy=0,\\
 z^*Hz&=-T,\qquad y^*Hy=-T',\qquad T'=T(z')>0.
\end{align*}
The Gram matrix of \([E,F]\) is
\(\left(\begin{smallmatrix}0&C\\C^*&0\end{smallmatrix}\right)\).
It is invertible, so these columns are independent.
Replacing \(F\) by \(FC^{-1}\) gives the form
\(\left(\begin{smallmatrix}0&I\\I&0\end{smallmatrix}\right)\),
of signature \((r,r)\). Thus \(U\) is nondegenerate and \(U^\perp\)
has signature \((k,1)\).

Substitution gives
\[
 E^*Hx=F^*Hx=E^*Hy_0=F^*Hy_0=0.
\]
For example,
\(F^*Hx=c^*-C^*(C^*)^{-1}c^*=0\) and
\(E^*Hy_0=b-CC^{-1}b=0\).
The other two equations use the displayed isotropy and support
relations. Hence \(x,y_0\in U^\perp\).
Each removed vector is isotropic and orthogonal to the original
vector, giving \(h(x,x)=-T\) and \(h(y_0,y_0)=-T'\).
Since \(((C^*)^{-1}c^*)^*=cC^{-1}\), expansion gives
\[
 h(x,y_0)=x^*Hy_0
 =d-cC^{-1}b-cC^{-1}b+cC^{-1}b=s.
\]
In the global conjugate-first convention,
\(h(x,p)=s+(s/T)(-T)=0\). Splitting off the negative line
\(\mathbb Cx\) from the space of signature \((k,1)\) proves that
\(V_+\) has signature \((k,0)\). Expanding the norm gives
\[
 h(p,p)=-T'+|s|^2/T.
\]
Subtracting \(cC^{-1}\) times the first \(r\) rows from the last
row of \(\widehat Z\) yields
\(D_{r+1}=(\det C)s\). Substitution in the actual definition of
\(G_f\) proves the boxed formula.
\end{proof}
For \(r=0\), the empty matrix is invertible with determinant \(1\),
\(U=0\), and this is the usual Lorentz projection.
For \(k=0\), \(V_+=0\) and the scalar vanishes on this chart;
globally the \(2r+2\) Gram vectors lie in dimension \(2r+1\), so
the source Gram identity gives \(G_f=0\) everywhere.
If \(D_r=0\), no inverse or limiting formula is used:
the definition gives \(G_f=|D_{r+1}|^2\ge0\) directly.
The assumption \(T>0\) is retained, as required by the source.

The determinant line and the positive space in this factorization
depend on the input configuration. Fixing the first vector at \(r=0\)
fixes a positive space in which the scalar is the norm of a projected
second vector; it does not make the two-variable scalar a Gram
inner product. Part (a) proves that such a reinterpretation of the
cross-values is impossible. Thus this calculation gives valid positive
geometry and a precise direct-kernel obstruction, without a historical
identification or a Xi approximation.
